% =============================================================================
% CataractAI: Perspektywa częstotliwościowa w kuracji zbiorów danych
% Dokumentacja techniczna (Polski)
% =============================================================================
\documentclass[11pt,a4paper]{article}

% --- Encoding & Language -----------------------------------------------------
\usepackage[utf8]{inputenc}
\usepackage[T1]{fontenc}
\usepackage[polish]{babel}

% --- Mathematics -------------------------------------------------------------
\usepackage{amsmath,amssymb,amsfonts}

% --- Tables ------------------------------------------------------------------
\usepackage{booktabs}
\usepackage{multirow}
\usepackage{array}

% --- Graphics ----------------------------------------------------------------
\usepackage{graphicx}

% --- Hyperlinks --------------------------------------------------------------
\usepackage[colorlinks=true,linkcolor=blue!70!black,citecolor=green!50!black,
            urlcolor=blue!60!black]{hyperref}

% --- Page Layout -------------------------------------------------------------
\usepackage[margin=2.5cm]{geometry}

% --- Headers -----------------------------------------------------------------
\usepackage{fancyhdr}
\pagestyle{fancy}
\fancyhf{}
\fancyhead[L]{\small CataractAI: Perspektywa częstotliwościowa}
\fancyhead[R]{\small\thepage}
\renewcommand{\headrulewidth}{0.4pt}

% --- Title -------------------------------------------------------------------
\title{\textbf{CataractAI: Perspektywa częstotliwościowa\\
       w kuracji zbiorów danych}\\[0.5em]
       \large Obciążenie spektralne, cechy częstotliwościowe\\
       i ich wpływ na trening DETR}
\author{Bartłomiej Piotrowski\\
        \small Politechnika Warszawska\\
        \small Projekt doktorski --- ViTParticleFilterTracker}
\date{Marzec 2026}

% =============================================================================
\begin{document}
\maketitle

\begin{abstract}
Niniejszy dokument przedstawia teoretyczne i empiryczne uzasadnienie
włączenia cech z dziedziny Fouriera do pipeline'u selekcji zbiorów danych
do treningu detektorow narzędzi chirurgicznych opartych na DETR.  Opieramy
się na trzech zbieżnych nurtach badawczych: (i)~obciążeniu spektralnym sieci
neuronowych, (ii)~ulepszeniach częstotliwościowych we współczesnych
architekturach DETR oraz (iii)~selekcji coresetów zorientowanej na pokrycie.
Walidacja empiryczna wykazuje, ze cechy Fouriera maja $3{,}4\times$ wyższą
ważność per-wymiar niż embeddingi DINO, a wybrany podzbiór 20\,000 obrazow
zachowuje rozkład spektralny pełnej puli 90\,389 obrazow (test KS,
wszystkie $p > 0{,}35$).
\end{abstract}

\tableofcontents
\newpage

% =============================================================================
% Main content
% =============================================================================
% =============================================================================
% Badania: Obciążenie spektralne, cechy częstotliwościowe i kuracja datasetu
% dla detekcji narzędzi chirurgicznych oparte na DETR
% =============================================================================
% Samodzielna sekcja do artykułu IEEE ACCESS
% Kompilacja: \usepackage{amsmath,amssymb,booktabs,multirow,graphicx}
% =============================================================================

\section{Perspektywa częstotliwościowa w kuracji zbiorów danych}
\label{sec:fourier_research}

Kompozycja danych treningowych w sposób fundamentalny determinuje, czego sieć
neuronowa jest w stanie się nauczyć.  W niniejszej sekcji przedstawiamy
teoretyczne i empiryczne uzasadnienie włączenia cech z dziedziny Fouriera do
pipeline'u selekcji zbiorów danych.  Opieramy się na trzech zbieżnych nurtach
badawczych: (i)~obciążeniu spektralnym sieci neuronowych, (ii)~ulepszeniach
częstotliwościowych we współczesnych architekturach DETR oraz
(iii)~selekcji coresetów zorientowanej na pokrycie.


% =============================================================================
\subsection{Obciążenie spektralne w głębokich sieciach neuronowych}
\label{sec:spectral_bias_theory}

\subsubsection{Zasada częstotliwościowa}

Rahaman i~wsp.~\cite{rahaman2019spectral} wykazali, ze głębokie sieci z
aktywacja ReLU wykazują systematyczne \emph{obciążenie spektralne} (ang.\
\emph{spectral bias}): składowe niskoczęstotliwościowe funkcji docelowej są
uczone jako pierwsze, podczas gdy składowe wysokoczęstotliwościowe zbiegaja
o rzędy wielkości wolniej.  Zjawisko to, znane również jako \emph{zasada
częstotliwościowa} (ang.\ \emph{frequency principle}, F-Principle), zostalo
sformalizowane za pomoca jądra neuronowego stycznego (Neural Tangent Kernel,
NTK)~\cite{jacot2018ntk}.

Niech $f_{\boldsymbol{\theta}}: \mathbb{R}^d \to \mathbb{R}$ oznacza sieć
o~parametrach $\boldsymbol{\theta}$, trenowana na danych
$\{(\mathbf{x}_i, y_i)\}_{i=1}^N$ metodą gradientu prostego ze współczynnikiem
uczenia $\eta$.  Jądro NTK definiujemy jako:

\begin{equation}
    \Theta(\mathbf{x}, \mathbf{x}')
    = \left\langle
        \frac{\partial f_{\boldsymbol{\theta}}(\mathbf{x})}
             {\partial \boldsymbol{\theta}},\;
        \frac{\partial f_{\boldsymbol{\theta}}(\mathbf{x}')}
             {\partial \boldsymbol{\theta}}
      \right\rangle
\end{equation}

W granicy nieskończonej szerokości jądro $\Theta$ pozostaje stałe podczas
treningu, a dynamika uczenia staje się liniowa.  Rozklad błędu rezydualnego
w bazie własnej $\Theta$ daje:

\begin{equation}
    e_k(t) = e_k(0) \cdot \exp(-\eta \lambda_k t)
    \label{eq:ntk_convergence}
\end{equation}

gdzie $\lambda_k$ jest $k$-ta wartością własną, a $e_k(t)$ -- rzutem błędu na
$k$-ta funkcje własną w kroku treningowym~$t$.  Dla standardowych sieci MLP
wartości własne zanikają szybko ze wzrostem częstotliwości:

\begin{equation}
    \lambda_k \propto \|\boldsymbol{\omega}_k\|^{-\beta},
    \qquad \beta > 0
    \label{eq:eigenvalue_decay}
\end{equation}

gdzie $\boldsymbol{\omega}_k$ jest czestotliwoscia powiązaną z~$k$-ta funkcja
własną.  Zanik ten oznacza, ze składowe błędu na częstotliwości
$\|\boldsymbol{\omega}\|$ zbiegaja jak
$\exp(-\eta \|\boldsymbol{\omega}\|^{-\beta} t)$ --- wykładniczo wolniej
dla wyższych częstotliwości.

\subsubsection{Potwierdzenie empiryczne i strategie mitygacji}

Najnowsze prace zarówno potwierdzily, jak i próbowały złagodzić obciążenie
spektralne:

\begin{itemize}
    \item \textbf{Regularyzacja częstotliwościowa}~\cite{freqreg2025}:
          Analiza empiryczna współczesnych sieci CNN ujawnia, ze regularyzacja
          L2 tłumi akumulację energii wysokoczęstotliwościowej ponad
          $3\times$ w porównaniu do modeli bez regularyzacji, działając jako
          niejawny filtr dolnoprzepustowy na uczone cechy.

    \item \textbf{Wielostopniowe uczenie głębokie}~\cite{mgdl2024}: Adresuje
          obciążenie spektralne przez kompozycje wielu płytkich sieci
          neuronowych (SNN), gdzie każdy stopien uchwytuje
          niskoczęstotliwościowe rezydua poprzedniego stopnia.  Ich złożenia
          mogą reprezentować funkcje wysokoczęstotliwościowe, które sieci
          monolityczne nie są w stanie odtworzyć.

    \item \textbf{Mapowanie cech Fouriera}~\cite{tancik2020fourier}:
          Przepuszczenie wejść przez $\gamma(\mathbf{x}) =
          [\cos(2\pi\mathbf{B}\mathbf{x}),
          \sin(2\pi\mathbf{B}\mathbf{x})]$ z losowymi częstotliwościami
          $\mathbf{B}$ przekształca jądro NTK w jądro stacjonarne z
          regulowana szerokością pasma, umożliwiając sieciom MLP uczenie
          funkcji wysokoczęstotliwościowych.
\end{itemize}

\subsubsection{Implikacje dla danych treningowych}

Obciążenie spektralne ma bezpośrednią konsekwencję dla kompozycji zbioru
danych.  Szybkość zbieżności na częstotliwości $\boldsymbol{\omega}$ zależy
nie tylko od wartości własnej NTK $\lambda(\boldsymbol{\omega})$, ale także
od \emph{energii spektralnej danych treningowych} na tej częstotliwości.
Rozważmy transformate Fouriera empirycznego rozkładu danych:

\begin{equation}
    \hat{p}_{\text{data}}(\boldsymbol{\omega})
    = \frac{1}{N} \sum_{i=1}^{N}
      \hat{I}_i(\boldsymbol{\omega})
\end{equation}

Jesli zbior treningowy niedoreprezentuje obrazow o silnej zawartości
wysokoczęstotliwościowej (tj.\
$\hat{p}_{\text{data}}(\boldsymbol{\omega}) \approx 0$ dla dużych
$\|\boldsymbol{\omega}\|$), to nawet z korekcjami NTK sieć otrzymuje
niewystarczający sygnał gradientu do uczenia się wzorców
wysokoczęstotliwościowych.  Motywuje to \emph{częstotliwościowo świadomą
kurację zbiorów danych}: selekcję podzbioru treningowego, który zachowuje
rozkład spektralny domeny docelowej.


% =============================================================================
\subsection{Dziedzina częstotliwościowa we współczesnych architekturach DETR}
\label{sec:detr_frequency}

Znaczenie przetwarzania częstotliwościowego dla detekcji obiektów zostalo
rozpoznane przez falę najnowszych wariantow DETR, które jawnie włączają
operacje spektralne do swoich architektur.
Tabela~\ref{tab:freq_detr_variants} podsumowuje te metody.

\begin{table*}[t]
\centering
\caption{Najnowsze warianty DETR z modułami wzmocnienia częstotliwościowego
         (2024--2025).  Wszystkie metody wykazują, ze zachowanie lub wzmocnienie
         cech wysokoczęstotliwościowych poprawia dokładność detekcji,
         szczególnie dla małych obiektów.}
\label{tab:freq_detr_variants}
\begin{tabular}{llccp{5.5cm}}
\toprule
\textbf{Metoda} & \textbf{Publikacja} &
    \textbf{$\Delta$mAP} & \textbf{Kluczowa metryka} &
    \textbf{Mechanizm częstotliwościowy} \\
\midrule
FDA-DETR~\cite{fda_detr2025}
    & PMC 2025
    & +4{,}0\,pp
    & AP małych obj.
    & Dekompozycja falkowa 2D (LL/LH/HL/HH) z uczalnymi
      wagami pasm $(\alpha, \beta, \gamma, \delta)$ \\[4pt]
DFIR-DETR~\cite{dfir_detr2025}
    & arXiv 2025
    & --
    & 92{,}9\% mAP50
    & Optymalizacja z ograniczeniami w dziedzinie spektralnej
      zachowującą składowe granic wysokoczęstotliwościowych \\[4pt]
Freq-DETR~\cite{freq_detr2025}
    & Expert Syst.\ 2025
    & +4{,}9\%
    & VisDrone mAP
    & Dwugałęziowe przetwarzanie przestrzenne + częstotliwościowe
      z modułem FECM \\[4pt]
FD-RTDETR~\cite{fd_rtdetr2025}
    & Electronics 2025
    & +1{,}9\,pp
    & COCO mAP
    & Interakcja cech wewnątrz-skalowa oparta na częstotliwości
      i uwadze (Frequency-Attention) \\
\bottomrule
\end{tabular}
\end{table*}

\subsubsection{Dekompozycja falkowa w FDA-DETR}

FDA-DETR~\cite{fda_detr2025} dokonuje dekompozycji map cech za pomoca
dyskretnej transformaty falkowej 2D:

\begin{equation}
    \psi_{j,m,n}(x,y) = 2^{j/2}\,\psi(2^j x - m,\; 2^j y - n)
\end{equation}

co daje cztery pod-pasma na każdej skali: $\mathbf{F}_{\text{LL}}$
(nisko-nisko, struktura globalna), $\mathbf{F}_{\text{LH}}$ (krawędzie
poziome), $\mathbf{F}_{\text{HL}}$ (krawędzie pionowe) oraz
$\mathbf{F}_{\text{HH}}$ (detale diagonalne).  Wzmocniona mapa cech:

\begin{equation}
    \mathbf{F}_{\text{enh}} = \alpha\,\mathbf{F}_{\text{LL}}
                            + \beta\,\mathbf{F}_{\text{LH}}
                            + \gamma\,\mathbf{F}_{\text{HL}}
                            + \delta\,\mathbf{F}_{\text{HH}}
\end{equation}

gdzie wagi $\alpha, \beta, \gamma, \delta$ są uczone oddzielnie dla każdej
skali.  Badania ablacyjne pokazują, ze stopniowe dodawanie składowych
wysokoczęstotliwościowych poprawia AP małych obiektów z~54{,}8\%
(tylko LL) do 57{,}5\% (wszystkie pod-pasma), co potwierdza, ze
\textbf{informacja o granicach w wysokich czestotliwosciach jest krytyczna
dla detekcji małych narzędzi chirurgicznych}.

\subsubsection{Optymalizacja spektralną w DFIR-DETR}

DFIR-DETR~\cite{dfir_detr2025} formalizuje agregacje cech jako optymalizacje
z ograniczeniami w dziedzinie spektralnej:

\begin{equation}
    \min_{\mathbf{F}}
    \|\mathcal{F}[\mathbf{F}] - \mathcal{F}[\mathbf{F}_{\text{target}}]\|^2
    + \lambda\,\Omega(\mathbf{F})
\end{equation}

gdzie $\mathcal{F}[\cdot]$ oznacza 2D FFT, a $\Omega$ jest regularyzatorem
penalizującym utrate składowych wysokoczęstotliwościowych.  Bezpośrednio
zachowuje to detale granic, które konwolucje przestrzenne stopniowo wygładzają.

\subsubsection{Znaczenie dla kuracji zbiorów danych}

Te innowacje architektoniczne ujawniają fundamentalną właściwość:
\textbf{jakość detekcji DETR jest ograniczona przez zawartość
wysokoczęstotliwościowa dostępną w mapach cech}.  Podczas gdy FDA-DETR
i~DFIR-DETR adresują to na poziomie architektury, ortogonalnym
i~komplementarnym podejściem jest zapewnienie, ze \emph{same dane treningowe}
zawierają wystarczającą różnorodność wysokoczęstotliwościowa.  Spektralnie
zubożone dane treningowe ograniczają cechy, które jakakolwiek architektura
może wyekstrahować, niezależnie od jej modułów wzmocnienia
częstotliwościowego.


% =============================================================================
\subsection{Selekcja coresetów zorientowana na pokrycie}
\label{sec:coreset_coverage}

Selekcja coresetów polega na znalezieniu małego podzbioru
$\mathcal{S} \subset \mathcal{D}$, który zachowuje właściwości uczenia
pełnego zbioru danych.  Najnowsze przeglądy i metody~\cite{coreset_survey2025}
identyfikują dwa komplementarne kryteria selekcji:

\begin{enumerate}
    \item \textbf{Selekcja oparta na trudności}: Priorytetyzacja próbek,
          które są trudne dla modelu (wysoka strata, wysoki wynik
          EL2N~\cite{paul2021deep}).  Sa one najbardziej informatywne dla
          uczenia granic decyzyjnych.

    \item \textbf{Selekcja oparta na pokryciu}: Zapewnienie, ze wybrany
          podzbiór obejmuje pełny rozkład rozmaitości danych.  Metody
          obejmuja $k$-Center Greedy (minimaksowa lokalizacja
          usług)~\cite{coreset_survey2025}, Herding (dopasowanie
          momentów)~\cite{welling2009herding} oraz lokalizację usługowa
          (submodularne pokrycie)~\cite{facility_location}.
\end{enumerate}

\subsubsection{Luka pokrycia w selekcji czysto semantycznej}

Większość metod coresetowych działa w pojedynczej przestrzeni embeddingowej
--- typowo w przedostatniej warstwie wytrenowanego klasyfikatora.  Blind
Coreset Selection~\cite{blindcs2025} używa embeddingów CLIP;
ELFS~\cite{elfs2025} używa cech z głębokiego klasteringu.  Te
\emph{semantyczne} embeddingi chwytają ,,co jest na obrazie'', ale są
w~znacznej mierze niezmiennicze względem statystyk niskopoziomowych, takich
jak ostrość, poziom szumu czy zawartość częstotliwościowa.

Rozważmy dwie klatki z operacji zacmy pokazujace ten sam czubek
fakoemulsyfikatora w podobnych pozycjach.  Embedding DINO lub CLIP odwzorowałby
je na bliskie punkty ($\|\mathbf{z}_a - \mathbf{z}_b\| \approx 0$), jednak
jedna może być w ostrej fokusie z wysokim $R_{\text{HL}}$ (stosunek energii
wysokich do niskich częstotliwości), a druga rozmyta z niskim $R_{\text{HL}}$.
Selekcja czysto semantyczna może odrzucić rozmyty wariant jako ,,redundantny'',
mimo ze dostarcza on cennego sygnału treningowego do obsługi ramek
nieostrych podczas inferencji.

\subsubsection{Cechy Fouriera jako komplementarną os pokrycia}

Dodanie cech Fouriera do kryterium selekcji wprowadza \emph{os pokrycia
w dziedzinie częstotliwościowej}, ortogonalna do pokrycia semantycznego:

\begin{equation}
    d^2(i,j) = \underbrace{\alpha_{\text{sem}}^2\,
                \|\hat{\mathbf{x}}_i^{\text{sem}} - \hat{\mathbf{x}}_j^{\text{sem}}\|^2}
               _{\text{odległość semantyczna (DINO)}}
             + \underbrace{\alpha_{\text{four}}^2\,
                \|\hat{\mathbf{x}}_i^{\text{four}} - \hat{\mathbf{x}}_j^{\text{four}}\|^2}
               _{\text{odległość spektralną (Fourier)}}
    \label{eq:combined_distance}
\end{equation}

Zapewnia to, ze obrazy różniące się zawartością częstotliwościowa, ale nie
semantyczna, są przypisywane do \emph{różnych} klastrów, zapobiegając
spektralnej jednorodności wybranego podzbioru.

Framework coresetowy InfoMax~\cite{infomax2025} dostarcza teoretycznego
wsparcia: optymalny coreset maksymalizuje informację wzajemną między
podzbiorem a pelnym zbiorem danych.  Ponieważ cechy Fouriera chwytają
informację ortogonalne do cech semantycznych
(Sekcja~\ref{sec:feature_importance}), ich włączenie scisle zwiększa
informację wzajemną selekcji, poprawiając trening.


% =============================================================================
\subsection{Analiza częstotliwościowa danych z wideo chirurgicznego}
\label{sec:surgical_fourier}

Wideo z operacji zacmy prezentuje unikalne charakterystyki spektralne, które
motywują częstotliwościowo świadomą kurację zbiorów danych.

\subsubsection{Zrodla zmienności częstotliwościowej}

\begin{enumerate}
    \item \textbf{Krawedzie narzędzi}: Czubki fakoemulsyfikatora, pincety
          i~injektory IOL tworzą ostre granice na tle miekkiej tkanki,
          generując silną zawartość wysokoczęstotliwościowa
          ($R_{\text{HL}} > 15$).

    \item \textbf{Odbicia spekularne}: Mokra powierzchnia rogówki generuje
          zlokalizowane szczyty wysokoczęstotliwościowe w widmie Fouriera,
          z~energia skoncentrowana w wąskich pasmach kierunkowych.

    \item \textbf{Tekstura tkanki}: Tęczówka, torebka soczewki i materiał
          korowy wykazują charakterystyczne tekstury średnioczęstotliwościowe
          ($0{,}1 \leq r < 0{,}5$), które różnią się w zależności od fazy
          operacji.

    \item \textbf{Zmienność ostrości}: Fokus mikroskopu zmienia się w trakcie
          operacji, tworząc ramki od ostrych
          ($\hat{H}_{\text{spec}} \approx 0{,}999$, wysoka entropia) do
          rozmytych ($\hat{H}_{\text{spec}} \approx 0{,}984$, niższa
          entropia).

    \item \textbf{Redundancja temporalna}: Przy 25--30\,kl./s kolejne klatki
          są niemal identyczne zarówno w dziedzinie przestrzennej, jak
          i~częstotliwościowej.  Filtrowanie podobienstwa Fouriera
          ($\tau_{\text{sim}} = 0{,}7$) efektywnie usuwa te
          prawie-duplikaty przed kosztowną ekstrakcją cech.
\end{enumerate}

\subsubsection{Empiryczna różnorodność spektralną}

Analiza wybranego zbioru danych ($|\mathcal{S}| = 20\,000$ obrazow) ujawnia
znaczącą różnorodność spektralną (Tabela~\ref{tab:spectral_stats}).

\begin{table}[h]
\centering
\caption{Statystyki spektralne wybranego podzbioru treningowego.
         Wysoki współczynnik zmienności (CV) $R_{\text{HL}}$ potwierdza
         zachowanie zarówno ostrych, jak i rozmytych próbek.}
\label{tab:spectral_stats}
\begin{tabular}{lccccc}
\toprule
\textbf{Metryka} & $\min$ & $\mu$ & $\tilde{x}$ & $\sigma$ & $\max$ \\
\midrule
$R_{\text{HL}}$
    & 4{,}79 & 8{,}47 & 6{,}91 & 5{,}35 & 28{,}79 \\
$\hat{H}_{\text{spec}}$
    & 0{,}984 & 0{,}989 & -- & -- & 0{,}999 \\
$f_c$
    & 0{,}416 & 0{,}442 & -- & -- & 0{,}533 \\
\bottomrule
\multicolumn{6}{l}{\small CV($R_{\text{HL}}$) $= \sigma/\mu = 63{,}2\%$} \\
\end{tabular}
\end{table}

$6$-krotny zakres $R_{\text{HL}}$ (od~4{,}79 do~28{,}79) obejmuje zarówno
próbki \emph{zdominowane przez gładkość} (rozmyte klatki, niska
szczegółowość), jak i próbki \emph{zdominowane przez częstotliwość} (ostre
krawędzie narzędzi, odbicia spekularne).
Tabela~\ref{tab:spectral_extremes} zestawia oba ekstrema.

\begin{table}[h]
\centering
\caption{Porównanie sygnatur spektralnych: próbki o najnizszym i najwyzszym
         $R_{\text{HL}}$ w wybranym zbiorze danych.}
\label{tab:spectral_extremes}
\begin{tabular}{lcc}
\toprule
& \textbf{Niski $R_{\text{HL}}$} & \textbf{Wysoki $R_{\text{HL}}$} \\
& (gładki, rozmyty) & (ostry, szczegółowy) \\
\midrule
$R_{\text{HL}}$             & 4{,}79  & 28{,}79 \\
$\hat{E}_{\text{low}}$      & 0{,}072 & 0{,}019 \\
$\hat{E}_{\text{mid}}$      & 0{,}582 & 0{,}439 \\
$\hat{E}_{\text{high}}$     & 0{,}346 & 0{,}542 \\
$\hat{H}_{\text{spec}}$     & 0{,}984 & 0{,}999 \\
$f_c$                       & 0{,}416 & 0{,}533 \\
Anizotropia                 & 1{,}353 & 1{,}440 \\
\bottomrule
\end{tabular}
\end{table}


% =============================================================================
\subsection{Walidacja empiryczna}
\label{sec:fourier_validation}

\subsubsection{Analiza ważności cech}
\label{sec:feature_importance}

W celu kwantyfikacji wkładu każdej grupy cech do jakości selekcji zbioru
danych wytrenowaliśmy regresor Random Forest~\cite{breiman2001random}
do przewidywania wyników trudności EL2N z pozostalych cech (DINO, Fourier,
SAM).  Wyniki ważności Giniego (Tabela~\ref{tab:rf_importance}) ujawniają,
ze cechy Fouriera są niemal równie informatywne jak embeddingi DINO do
przewidywania trudności próbek.

\begin{table}[h]
\centering
\caption{Waznosc cech Random Forest do przewidywania trudności EL2N.
         Waznosc per-wymiar obliczono jako
         $\text{Waznosc}_g / d_g$.}
\label{tab:rf_importance}
\begin{tabular}{lcccc}
\toprule
\textbf{Cecha} & $d_g$ & \textbf{Waznosc} &
    \textbf{Per-wym.} & \textbf{Waga} \\
\midrule
DINO (PCA)       & 32 & 48{,}4\% & 1{,}51\% & 0{,}35 \\
\textbf{Fourier} &  9 & \textbf{46{,}3\%} & \textbf{5{,}14\%} & 0{,}15 \\
SAM              &  1 &  5{,}3\% & 5{,}30\% & 0{,}20 \\
\bottomrule
\end{tabular}
\end{table}

\paragraph{Kluczowa obserwacja.}
Cechy Fouriera wykazują $5{,}14 / 1{,}51 = 3{,}4\times$ wyższą ważność
per-wymiar niż embeddingi DINO.  Ta nieproporcjonalną informatywność wskazuje,
ze statystyki dziedziny częstotliwościowej chwytają komplementarną os
zmienności --- taka, której wysokopoziomowe reprezentacje semantyczne DINO
nie kodują.  Korelacje te można wytłumaczyć mechanizmem obciążenia
spektralnego: obrazy, które są ,,trudne'' dla DETR (wysoki EL2N), typowo
zawierają albo (a)~atypową zawartość wysokoczęstotliwościowa (ostre odbicia,
nietypowe tekstury), albo (b)~atypowe wzorce niskoczęstotliwościowe (brak
ostroci, zasłonięte narzędzia) --- oba typy są dobrze chwytane przez
9-wymiarowy wektor Fouriera.

\subsubsection{Zachowanie rozkładu}

Warunkiem koniecznym prawidłowej selekcji zbioru danych jest zachowanie przez
wybrany podzbiór $\mathcal{S}$ właściwości statystycznych pełnej puli
$\mathcal{D}$.  Weryfikujemy to za pomoca dwu-próbkowego testu
Kolmogorowa--Smirnowa~(KS):

\begin{equation}
    D_{\text{KS}} = \sup_x \bigl|F_{\mathcal{D}}(x) - F_{\mathcal{S}}(x)\bigr|
\end{equation}

gdzie $F_{\mathcal{D}}$ i $F_{\mathcal{S}}$ są empirycznymi dystrybuantami
danej metryki nad pulą ($N = 90\,389$) i podzbiorem
($|\mathcal{S}| = 20\,000$).

\begin{table}[h]
\centering
\caption{Test Kolmogorowa--Smirnowa dla zachowania cech Fouriera.
         Wszystkie wartości $p > 0{,}05$: wybrany podzbiór jest statystycznie
         nieodróżnialny od pełnej puli w dziedzinie częstotliwościowej pomimo
         $4{,}5\times$ redukcji rozmiaru.}
\label{tab:ks_fourier}
\begin{tabular}{lccc}
\toprule
\textbf{Metryka Fouriera} &
    $\mu_{\mathcal{D}}$ & $\mu_{\mathcal{S}}$ & wartosc $p$ \\
\midrule
$\hat{E}_{\text{low}}$      & 0{,}0525 & 0{,}0524 & 0{,}506 \\
$\hat{E}_{\text{mid}}$      & 0{,}5704 & 0{,}5706 & 0{,}820 \\
$\hat{E}_{\text{high}}$     & 0{,}3769 & 0{,}3770 & 0{,}928 \\
$\hat{H}_{\text{spec}}$     & 0{,}9894 & 0{,}9895 & 0{,}410 \\
$f_c$                       & 0{,}4415 & 0{,}4416 & 0{,}356 \\
$R_{\text{HL}}$             & 8{,}468  & 8{,}471  & 0{,}991 \\
\bottomrule
\end{tabular}
\end{table}

Wszystkie wartości $p$ znaczaco przekraczają poziom istotności
$\alpha = 0{,}05$, potwierdzając, ze procedura selekcji zachowuje wierność
spektralną pomimo redukcji zbioru danych do 22{,}1\% jego pierwotnego
rozmiaru.


% =============================================================================
\subsection{Mechanizm dzialania: jak cechy Fouriera poprawiaja trening DETR}
\label{sec:mechanism}

Integracja cech Fouriera w kryterium selekcji wpływa na trening DETR
poprzez trzy komplementarne mechanizmy.

\subsubsection{Pokrycie spektralne}

Wybierając próbki obejmujące pełny zakres $R_{\text{HL}}$ (4{,}79--28{,}79),
zbior treningowy eksponuje DETR zarówno na niskoczęstotliwościowe tla
(gładką tkanka), jak i na wysokoczęstotliwościowe obiekty pierwszego planu
(ostre krawędzie narzędzi).  Zgodnie z Równ.~\ref{eq:ntk_convergence},
szybkość zbieżności dla wysokoczęstotliwościowych składowych błędu wynosi:

\begin{equation}
    e_{\text{HF}}(t) \propto \exp\!\bigl(-\eta \lambda_{\text{HF}} t\bigr)
    \cdot \underbrace{\hat{p}_{\text{data}}(\boldsymbol{\omega}_{\text{HF}})}
          _{\text{energia spektralną danych}}
\end{equation}

Spektralnie zróżnicowany zbior treningowy zwiększa
$\hat{p}_{\text{data}}(\boldsymbol{\omega}_{\text{HF}})$, przyspieszając
zbieznosc cech wysokoczęstotliwościowych krytycznych dla lokalizacji granic.
Potwierdzają to badania ablacyjne FDA-DETR~\cite{fda_detr2025}: dodanie
wysokoczęstotliwościowych składowych falkowych (LH, HL, HH) poprawia AP
małych obiektów o~+2{,}7\,pp, demonstrując bezpośrednią wartosc informacji
wysokoczęstotliwościowej dla detekcji.

\subsubsection{Eliminacja redundancji}

Wideo chirurgiczne przy 25--30\,kl./s generuje masywną redundancję temporalna.
Wstepne filtrowanie podobienstwa Fouriera usuwa 10--20\% prawie-duplikatow
ramek \emph{przed} kosztowną ekstrakcją cech na GPU (DINO, EL2N):

\begin{equation}
    \text{sim}(i,j) = \frac{1}{1 + \|\hat{\mathbf{x}}_i^{\text{four}}
                      - \hat{\mathbf{x}}_j^{\text{four}}\|_2}
\end{equation}

Służy to podwójnemu celowi: (a)~redukcji kosztu obliczeniowego o~10--20\%,
oraz (b)~zapobieganiu dominacji klasteringu $k$-średnich przez redundantne
warunki obrazowania, co zmniejszyłoby efektywną różnorodność wybranego
podzbioru.

\subsubsection{Spektralne balansowanie uwzgledniajace trudność}

Wysoka korelacja między cechami Fouriera a trudnością EL2N (ważność 46{,}3\%,
Tabela~\ref{tab:rf_importance}) tworzy naturalną synergie: obrazy
o~nietypowych sygnaturach spektralnych maja tendencję do bycia trudniejszymi
dla DETR i na odwrot.  W sformułowaniu ważonego $k$-średnich
(Równ.~\ref{eq:combined_distance}), te spektralnie nietypowe próbki zajmują
odrębne regiony przestrzeni cech, tworząc własne klastry.  Ponieważ z każdego
klastra wybierany jest jeden reprezentant, trudne próbki o rzadkich profilach
spektralnych są naturalnie priorytetyzowane bez jawnego ważenia opartego na
trudności.


% =============================================================================
\subsection{Powiazania z pracami pokrewnymi}
\label{sec:fourier_connections}

Tabela~\ref{tab:related_work} pozycjonuje nasze podejście w szerszym
krajobrazie metod częstotliwościowych dla kuracji danych treningowych
i detekcji obiektów.

\begin{table*}[t]
\centering
\caption{Nasze podejście w kontekście: metody częstotliwościowe dla kuracji
         zbiorów danych i detekcji obiektów.}
\label{tab:related_work}
\begin{tabular}{p{3.2cm}p{2.5cm}p{3cm}p{5.5cm}}
\toprule
\textbf{Kategoria} & \textbf{Praca} & \textbf{Publikacja} &
    \textbf{Relacja do naszego podejścia} \\
\midrule
\multicolumn{4}{l}{\textit{Teoria obciążenia spektralnego}} \\
\quad Zasada częst.
    & Rahaman i~wsp.
    & ICML 2019
    & Motywuje częstotliwościowo zróżnicowane zbiory treningowe
      do przeciwdzialania wrodzonemu obciążeniu niskoczest. \\[3pt]
\quad Cechy Fouriera
    & Tancik i~wsp.
    & NeurIPS 2020
    & Mapuje wejscia na baze Fouriera aby dostroisc szerokość pasma
      NTK; my stosujemy ta sama zasade do \emph{selekcji danych},
      a nie architektury \\[3pt]
\quad Regul.\ częst.
    & Zheng i~wsp.
    & arXiv 2025
    & Wykazuje, ze regularyzacja L2 tłumi WCz o~$3\times$; wzmacnia
      potrzebę danych bogatych w~WCz \\[3pt]
\midrule
\multicolumn{4}{l}{\textit{Detekcja wzmocniona częstotliwościowo}} \\
\quad FDA-DETR
    & Li i~wsp.
    & PMC 2025
    & Wzmocnienie falkowe daje +4{,}0\,pp AP małych obj.;
      potwierdza, ze WCz to waskie gardlo \\[3pt]
\quad DFIR-DETR
    & Wang i~wsp.
    & arXiv 2025
    & Optymalizacja spektralną zachowuje granice; my zachowujemy
      granice \emph{wybierając próbki zróżnicowane w WCz} \\[3pt]
\midrule
\multicolumn{4}{l}{\textit{Selekcja coresetów}} \\
\quad Blind CS
    & Anonimowi
    & ICLR 2025
    & Pokrycie przestrzeni embeddingowej dla danych bez etykiet; my
      rozszerzamy pokrycie o embeddingi częstotliwościowe \\[3pt]
\quad Coreset dla OD
    & Lee i~wsp.
    & CVPR 2024W
    & Pierwsza metodą coresetowa dla detekcji obiektów; używa cech
      specyficznych dla detekcji, ale bez cech częstotliwościowych \\[3pt]
\quad EL2N / Data Diet
    & Paul i~wsp.
    & NeurIPS 2021
    & Przycinanie oparte na trudności; my łączymy z~pokryciem
      częstotliwościowym dla zbalansowanej selekcji \\[3pt]
\midrule
\multicolumn{4}{l}{\textit{Obrazowanie medyczne i dane chirurgiczne}} \\
\quad Cataract-1K
    & Ghamsarian i~wsp.
    & Nat.\ Sci.\ Data 2024
    & Największy zbior danych z operacji zacmy (1000 wideo); my
      rozszerzamy o częstotliwościowo świadomą selekcję podzbiorów \\[3pt]
\quad Augm.\ Fouriera
    & Xu i~wsp.
    & Pattern Recogn.\ 2023
    & Modyfikacja widma amplitudowego do generalizacji między-domenowej;
      potwierdza: faza = semantyka, amplituda = statystyki
      niskopoziomowe \\
\bottomrule
\end{tabular}
\end{table*}

\paragraph{Nowość naszego podejścia.}
Podczas gdy przetwarzanie częstotliwościowe bylo badane zarówno na poziomie
architektury (FDA-DETR, DFIR-DETR), jak i na poziomie augmentacji
(augmentacja oparta na Fourierze), zastosowanie cech z dziedziny Fouriera
do \emph{kuracji zbiorów danych} w detekcji obiektów pozostaje niezbadane.
Nasza praca wypełnia te lukę poprzez:
\begin{enumerate}
    \item Ekstrakcję kompaktowej 9-wymiarowej sygnatury Fouriera na obraz,
          chwytającej różnorodność spektralną, kierunkowość i entropie;
    \item Użycie tych cech zarówno do filtrowania redundancji, jak i do
          częstotliwościowo świadomego klasteringu obok cech semantycznych
          (DINO) i opartych na trudności (EL2N);
    \item Empiryczne wykazanie, ze cechy Fouriera maja $3{,}4\times$ wyższą
          ważność per-wymiar niż DINO do przewidywania trudności próbek,
          podczas gdy wybrany podzbiór zachowuje rozkład spektralny pełnego
          zbioru danych (test KS, wszystkie $p > 0{,}35$).
\end{enumerate}


% =============================================================================
% WPISY BIBLIOGRAFICZNE (dodac do glownego pliku .bib)
% =============================================================================
%
% @inproceedings{rahaman2019spectral,
%   title     = {On the Spectral Bias of Neural Networks},
%   author    = {Rahaman, Nasim and Barber, Aristide and Klartag, Boris and
%                Arjovsky, Martin and Bengio, Yoshua},
%   booktitle = {ICML},
%   year      = {2019}
% }
%
% @inproceedings{jacot2018ntk,
%   title     = {Neural Tangent Kernel: Convergence and Generalization
%                in Neural Networks},
%   author    = {Jacot, Arthur and Gabriel, Franck and Hongler, Cl{\'e}ment},
%   booktitle = {NeurIPS},
%   year      = {2018}
% }
%
% @inproceedings{tancik2020fourier,
%   title     = {Fourier Features Let Networks Learn High Frequency Functions
%                in Low Dimensional Domains},
%   author    = {Tancik, Matthew and Srinivasan, Pratul P. and Mildenhall, Ben
%                and Fridovich-Keil, Sara and Raghavan, Nithin and
%                Singhal, Utkarsh and Ramamoorthi, Ravi and Barron, Jonathan T.
%                and Ng, Ren},
%   booktitle = {NeurIPS},
%   year      = {2020}
% }
%
% @article{freqreg2025,
%   title   = {Frequency Regularization: Unveiling the Spectral Inductive Bias
%              of Deep Neural Networks},
%   author  = {Zheng, Yifei i inni},
%   journal = {arXiv:2512.22192},
%   year    = {2025}
% }
%
% @inproceedings{mgdl2024,
%   title     = {Addressing Spectral Bias of Deep Neural Networks by
%                Multi-Grade Deep Learning},
%   booktitle = {NeurIPS},
%   year      = {2024}
% }
%
% @inproceedings{paul2021deep,
%   title     = {Deep Learning on a Data Diet: Finding Important Examples
%                Early in Training},
%   author    = {Paul, Mansheej and Ganguli, Surya and Dziugaite, Gintare K.},
%   booktitle = {NeurIPS},
%   year      = {2021}
% }
%
% @article{breiman2001random,
%   title   = {Random Forests},
%   author  = {Breiman, Leo},
%   journal = {Machine Learning},
%   volume  = {45},
%   pages   = {5--32},
%   year    = {2001}
% }
%
% @inproceedings{blindcs2025,
%   title     = {Blind Coreset Selection: Efficient Pruning for Unlabeled Data},
%   booktitle = {ICLR},
%   year      = {2025}
% }
%
% @inproceedings{elfs2025,
%   title     = {{ELFS}: Label-Free Coreset Selection with Proxy Training
%                Dynamics},
%   booktitle = {ICLR},
%   year      = {2025}
% }
%
% @inproceedings{infomax2025,
%   title     = {{InfoMax} Coreset Selection},
%   booktitle = {ICLR},
%   year      = {2025}
% }
%
% @article{coreset_survey2025,
%   title   = {A Coreset Selection of Coreset Selection Literature:
%              Introduction and Recent Advances},
%   journal = {arXiv:2505.17799},
%   year    = {2025}
% }
%
% @article{fda_detr2025,
%   title   = {{FDA-DETR}: A Frequency-Aware {DETR} with Dynamic Query and
%              Adaptive Multi-Task Optimization for Oriented Small Object
%              Detection},
%   author  = {Li, H. i inni},
%   journal = {PMC},
%   year    = {2025}
% }
%
% @article{dfir_detr2025,
%   title   = {{DFIR-DETR}: Frequency Domain Enhancement and Dynamic Feature
%              Aggregation for Cross-Scene Small Object Detection},
%   author  = {Wang, Y. i inni},
%   journal = {arXiv:2512.07078},
%   year    = {2025}
% }
%
% @article{freq_detr2025,
%   title   = {{Freq-DETR}: Frequency-Aware Transformer for Real-Time Small
%              Object Detection in Unmanned Aerial Vehicle Imagery},
%   journal = {Expert Systems with Applications},
%   year    = {2025}
% }
%
% @article{fd_rtdetr2025,
%   title   = {{FD-RTDETR}: Frequency Enhancement and Dynamic Sequence-Feature
%              Optimization for Object Detection},
%   journal = {Electronics},
%   volume  = {14},
%   number  = {23},
%   pages   = {4715},
%   year    = {2025}
% }
%
% @article{cataract1k2024,
%   title   = {Cataract-1K Dataset for Deep-Learning-Assisted Analysis of
%              Cataract Surgery Videos},
%   author  = {Ghamsarian, Negin i inni},
%   journal = {Nature Scientific Data},
%   year    = {2024}
% }
%
% @article{fourier_augmentation2023,
%   title   = {Fourier-Based Augmentation with Applications to Domain
%              Generalization},
%   author  = {Xu, Q. i inni},
%   journal = {Pattern Recognition},
%   year    = {2023}
% }
%
% @inproceedings{dzanic2020fourier,
%   title     = {Fourier Spectrum Discrepancies in Deep Network Generated
%                Images},
%   author    = {Dzanic, Tarik and Shah, Karan and Witherden, Freddie},
%   booktitle = {NeurIPS},
%   year      = {2020}
% }
%
% @inproceedings{coreset_od2024,
%   title     = {Coreset Selection for Object Detection},
%   author    = {Lee, Hojun and Kim, Suyoung and Lee, Junhoo},
%   booktitle = {CVPR Workshop (DDCV)},
%   year      = {2024}
% }
%
% @inproceedings{welling2009herding,
%   title     = {Herding Dynamical Weights to Learn},
%   author    = {Welling, Max},
%   booktitle = {ICML},
%   year      = {2009}
% }


% =============================================================================
% Bibliography
% =============================================================================
\newpage
\bibliographystyle{IEEEtran}
\bibliography{references}

\end{document}
